\subsection{Tube type targets}
We prove Proposition~C, namely that whenever $G_{\mathbb R}$ has tube type and $n\geqslant2$, the Toledo invariant of a representation from $\Gamma$ to $G_{\mathbb R}$ satisfies an inequality stronger than the Milnor-Wood inequality, preventing a representation in such a group to be maximal.
\begin{proposition}
Let $\Gamma$ be a uniform lattice in $\mathrm{SU}(1,n)$, and let $X=\Gamma\backslash\mathbb H^n_{\mathbb C}$. Assume that the simple real algebraic Hermitian Lie group $G_{\mathbb R}$ has tube type and let $r_M$ be the real rank of $G_{\mathbb R}$. Let $\rho$ be a representation $\Gamma\to G_{\mathbb R}$. Then
\[
|\tau(\rho)|\leqslant\max\Big\{r_M-1,\frac{r_M}{2}\cdot\frac{n+1}{n}\Big\}\operatorname{vol}(X).
\]
\end{proposition}
\begin{remark}
Observe that $\max\big\{r_M-1,\frac{r_M}{2}\cdot\frac{n+1}{n}\big\}=r_M-1$ unless $r_M\leqslant2$ or $r_M=3$ and $n=2$.
\end{remark}
\begin{proof}
We may assume that $\Gamma$ is torsion free, that $\tau(\rho)>0$, that $\rho$ is reductive (see the beginning of Section~4), and that $G$ is simply connected (see Remarks~5.14 and~5.15). Then, the constructions of Sections~2,~3 and~4 are valid and the inequality of the proposition is equivalent to the inequality
\[
\deg_{\jepPubScript{T}}(E_0)\leqslant\max\Big(\frac{r_M-1}{2},\frac{r_M}{2}\cdot\frac{n+1}{2n}\Big)\deg_{\jepPubScript{T}}(L^\star).
\]
We use freely the notation of Section~4. If the generic rank of $\beta$ on the projectivized tangent bundle $\mathbb PT_X$ of $X$ is $\leqslant r_M-1$ then we are done by Theorem~4.6. Therefore we may assume that the generic rank of $\beta$ on $\mathbb PT_X$ is $r_M$.

We come back to the Higgs bundle $(\overline E,\overline\theta)$ on $X$ and we consider $\overline\beta:\overline E\otimes T_X\to\overline E$. The fact that $\operatorname{rk}\beta=r_M$ implies that $\overline\beta^{r_M}$, seen as a morphism from $\overline E_0\otimes\overline E_{r_M}^{\star}$ to the $r_M$-th symmetric power $S^{r_M}\Omega_X^1$ of $\Omega_X^1$, is not zero. Since $\Omega_X^1$ is a semistable bundle over $X$ ($X$ is Kähler-Einstein), so is $S^{r_M}\Omega_X^1$. On the other hand, $\overline E_0\otimes\overline E_{r_M}^{\star}$ is also semistable because it is a line bundle by Section~2.4.4. Therefore $\mu_\omega(\overline E_0\otimes\overline E_{r_M}^{\star})\leqslant\mu_\omega(S^{r_M}\Omega_X^1)$ (recall that $\omega$ is our Kähler form on $X$), so that $\deg_\omega\overline E_0-\deg_\omega\overline E_{r_M}\leqslant r_M\mu_\omega(\Omega_X^1)$. Now, as explained in Section~2.4.4, the $K$-modules $\jepPubBbb{E}_{r_M}$ and $\jepPubBbb{E}_0^{\star}$ are isomorphic, so that $\deg_\omega\overline E_{r_M}=-\deg_\omega\overline E_0$ by Fact~4.5. We get the result, since by Equations~(4.2) and~(4.3) in Section~4.2 and the isomorphism $K_X\simeq\det\Omega_X^1$, we have $\deg_\omega(\Omega_X^1)=\frac{n+1}{2}\deg_{\jepPubScript{T}}(L^\star)$.
\end{proof}

{\footnotesize
